\section{Rigid motions and congruence of finite figures}
\label{sec:motions}
Let $J(x,y)=(-y,x)$ and write $[u,v]=u_1v_2-u_2v_1$.
Every distance-preserving map of the plane has the form
$X\mapsto t+UX$ for an orthogonal $2\times2$ matrix $U$. To prove this,
subtract the image of the origin, and call the resulting map $f$. The
vectors $u=f(1,0)$ and $v=f(0,1)$ have squared length one and squared
distance two, hence $u\cdot v=0$. For $X=(a,b)$, subtraction of its
squared distances to the origin and to $(1,0)$ gives
$f(X)\cdot u=a$. Comparison with $(0,1)$ similarly gives
$f(X)\cdot v=b$. Since $u,v$ form an orthonormal basis, $f(X)=au+bv$.
Conversely every displayed map preserves squared distances, has inverse
$Y\mapsto U^{\mathsf T}(Y-t)$, and preserves betweenness. Thus motions
can be represented by six scalar entries subject to $U^{\mathsf T}U=I$.
Assertions about such motions are first-order field assertions. The
argument about an arbitrary map is a schematic statement about a map
supplied in the metatheory; the finite code suffices inside arithmetic.

\begin{proposition}
Two labelled lists $P_1,\ldots,P_n$ and $Q_1,\ldots,Q_n$ are congruent by
a rigid motion, allowing reflection, if and only if all corresponding
pairwise squared distances agree. There is a PR procedure computing a
motion or returning a mismatched pair.
\end{proposition}
\begin{proof}
Compare all pairs. Necessity follows from the definition of motion.
For the converse, return the identity for empty lists. If all $P_i$
coincide, all $Q_i$ coincide, and translation by $Q_1-P_1$ works.
Otherwise choose a distinct pair, renumbered as $1,2$. Put
$u=P_2-P_1$, $v=Q_2-Q_1$, and $N=u\cdot u=v\cdot v>0$.
Write
\[
   P_i=P_1+a_i u+b_i Ju,\qquad Q_i=Q_1+c_i v+d_i Jv .
\]
The coefficients are dot products with the basis vectors divided by
$N$. Distances to points $1,2$ give $a_i=c_i$ by subtraction and then
$b_i^2=d_i^2$ from the distance to point $1$. If every $b_i=0$, every
$d_i=0$ and take $e=1$. Otherwise choose $k$ with $b_k\ne0$. Set
$e=d_k/b_k$, so $e^2=1$. Expanding the distance between points $i,k$
and using the known squares gives $b_i b_k=d_i d_k$. Dividing by $b_k$
and multiplying by $e$ yields $d_i=e b_i$ for all $i$. In both cases a
witness is the map
\[
 P_1+au+bJu\ \longmapsto\ Q_1+av+e bJv .
\]
The basis pairs are orthogonal and have common squared length $N$, so
this map is a rigid motion. Its determinant is $e$. In the collinear
case either sign works. In the noncollinear case the two basis
coefficients at $k$ force the sign, so this also decides congruence by
orientation-preserving motions. Every search ranges over a given finite
list.
\end{proof}

The pure point language records equidistance and leaves the unit of
length free. Its automorphisms on the real algebraic plane are exactly
the similarities, including those reversing orientation. Indeed, an
automorphism sends any frame to a frame, and the relational definitions
of Section~\ref{sec:frame-inverse} make its induced map on numerical
coordinates an ordered-field automorphism. Such an automorphism fixes
each rational and thus fixes each real algebraic number. The unique
zero of its record polynomial in its rational isolating interval
must go to the same zero. Consequently the point map is exactly the
frame change $O+au+bv\mapsto O'+au'+bv'$. Its linear part multiplies
all squared distances by the common positive factor
$(u'\cdot u')/(u\cdot u)$. Conversely these similarities preserve
betweenness and congruence, as follows directly from those formulas.

For unlabelled finite configurations of the same cardinality, enumerate the
permutations of the second list and apply the labelled procedure in order.
The first success gives both a matching and a witnessing motion. Conversely
any congruence of the configurations induces a permutation on the listed
points (with multiplicities matched if the lists describe multisets).
Exhausting the $n!$ possibilities therefore certifies failure.
